\subsection{EXPONENTIAL IN COMPLETE NORMED ASSOCIATIVE ALGEBRAS}

In this no. A will denote a unital associative algebra with a norm \(x \mapsto \| x \|\) satisfying the conditions:

\[
\begin{array}{c} \| x + y \| \leqslant \sup (\| x \|, \| y \|) \\ \| x y \| \leqslant \| x \|. \| y \| \\ \| 1 \| = 1 \end{array}
\]

for \(x, y\) in A, and complete with this norm. The results of the second and third paragraphs of § 7, no. 3, remain valid.

We take \(\mathbf{I} = \{\mathrm{U}\}\) and consider the images \(\tilde e\) and \(\tilde{l}\) of the series \(e(\mathrm{U}) = \sum_{n\geqslant 1}\frac{\mathrm{U}^n}{n!}\) and \(l(\mathrm{U}) = \sum_{n\geqslant 1}(-1)^{n - 1}\frac{U^n}{n}\) in \(\hat{\mathrm{P}} (\mathrm{A};\mathrm{A})\) . Then:

\begin{enumerate}
\def\labelenumi{(\arabic{enumi})}
\setcounter{enumi}{19}
\tightlist
\item
\end{enumerate}

\[
\left\| \left(\frac {\mathrm{U} ^ {n}}{n !}\right) ^ {\sim} \right\| \leqslant a ^ {- (n - 1) \theta}\tag{20}
\]

\[
\left\| \left(\frac {\mathrm{U} ^ {n}}{n}\right) ^ {\sim} \right\| \leqslant a ^ {- \frac {\log n}{\log p}}\tag{21}
\]

by Lemma 2 of no. 1. Hence the radius of absolute convergence of the series \(\tilde{e}\) (resp. \(\tilde{l}\) ) is \(\geqslant a^\theta\) (resp. \(\geqslant 1\) ) (Differentiable and Analytic Manifolds, R, 4.1.3). For \(R > 0\) , let \(G_R = \{x \in A \mid \|x\|<R\}\) ; we write \(G = G_\theta\) . The series \(\tilde{e}\) (resp. \(\tilde{l}\) ) defines an analytic mapping \(e_{A}\) (resp. \(l_{A}\) ) of \(G\) (resp. \(G_{1}\) ) into A. We write:

\begin{enumerate}
\def\labelenumi{(\arabic{enumi})}
\setcounter{enumi}{21}
\tightlist
\item
\end{enumerate}

\[
\exp_ {\mathrm{A}} (x) = 1 + e _ {\mathrm{A}} (x) = \sum_ {n \geqslant 0} \frac {x ^ {n}}{n !}\tag{22}
\]

\[
\log_ {\mathrm{A}} (x) = l _ {\mathrm{A}} (x - 1) = \sum_ {n \geqslant 1} (- 1) ^ {n - 1} \frac {(x - 1) ^ {n}}{n} \quad \text { for } x - 1 \in \dot {\mathrm{G}} _ {1}\tag{23}
\]

(we omit the index A when no confusion can arise). For \(x \in G_R\) and \(n \geqslant 1\) ,

\[
\left\| \frac {x ^ {n}}{n} \right\| \leqslant \left\| \frac {x ^ {n}}{n !} \right\| <   R ^ {n} a ^ {- (n - 1) \theta} = R \left(\frac {R}{a ^ {\theta}}\right) ^ {n - 1}\tag{24}
\]

and hence \(e_{\mathrm{A}}(\mathrm{G}_{\mathrm{R}}) \subset \mathrm{G}_{\mathrm{R}}, l_{\mathrm{A}}(\mathrm{G}_{\mathrm{R}}) \subset \mathrm{G}_{\mathrm{R}}\) for \(R \leqslant a^{\theta}\) .

PROPOSITION 4. Let R be a real number such that \(0 < R \leqslant a^{\theta}\) . The mapping \(\exp_{\mathbf{A}}\) defines an analytic isomorphism of \(\mathbf{G}_{R}\) onto \(1 + \mathbf{G}_{R}\) and the inverse isomorphism is the restriction of \(\log_{\mathbf{A}}\) to \(1 + \mathbf{G}_{R}\) .

\(e(l(\mathbf{X})) = l(e(\mathbf{X})) = \mathbf{X}\) . By (20), (21) and Differentiable and Analytic Manifolds, R, 4.1.5, we deduce that \(e_{\mathrm{A}}(l_{\mathrm{A}}(x)) = l_{\mathrm{A}}(e_{\mathrm{A}}(x))\) for \(x \in \mathrm{G}_{\mathrm{R}}\) . Then

\[
\begin{array}{l l} \exp_ {\mathrm{A}} (\log_ {\mathrm{A}} x) = x & \text {for} x \in 1 + \mathrm{G} _ {\mathrm{R}} \\ \log_ {\mathrm{A}} (\exp_ {\mathrm{A}} x) = x & \text {for} x \in \mathrm{G} _ {\mathrm{R}} \end{array}
\]

which completes the proof.

If A is given the bracket \([x,y]=xy-yx\) , A becomes a complete normed Lie algebra, for \(\|xy-yx\|\leqslant\sup(\|xy\|,\|yx\|)\leqslant\|x\|. \|y\|\) . Proposition 2 of no. 3 implies that the domain of absolute convergence of \(\tilde{H}\) contains \(G\times G\) and \(\tilde{H}\) therefore defines an analytic function \(h:G\times G\to A\) ; then

\[
h (x, y) = \sum_ {r, s \geqslant 0} \mathrm{H} _ {r, s} (x, y).\tag{25}
\]

PROPOSITION 5. For x, y in G,

\[
\exp_ {\mathrm{A}} x. \exp_ {\mathrm{A}} y = \exp_ {\mathrm{A}} h (x, y).\tag{26}
\]

\(e^{\mathrm{U}}e^{\mathrm{V}} = e^{\mathrm{H(U,V)}}\) and hence

\[
m \circ (1 + \tilde {e}, 1 + \tilde {e}) = (1 + \tilde {e}) \circ \tilde {\mathrm{H}}
\]

in \(\hat{\mathrm{P}}(\mathbf{A} \times \mathbf{A}; \mathbf{A})\) (where m denotes multiplication on A). The proposition then follows from Proposition 2, Lemma 3 and Differentiable and Analytic Manifolds, R, 4.1.5.


